%==============================================================================
% Appendix
%==============================================================================

\section{Score Definitions Reference}
\label{app:scores}

Table~\ref{tab:all-scores} lists all 16 score definitions registered with Langfuse.

\begin{table}[h]
\centering
\small
\caption{Complete score definitions. All use the 5-point categorical scale unless noted.}
\label{tab:all-scores}
\begin{tabular}{llp{7cm}}
\toprule
\textbf{Prefix} & \textbf{Name} & \textbf{Description} \\
\midrule
\multirow{5}{*}{threat\_} & overall\_quality & Holistic assessment of generated threats \\
& relevance\_to\_context & Match to application context \\
& completeness & Coverage of threat categories \\
& technical\_accuracy & Technical correctness \\
& hallucination\_score & Absence of fabricated content \\
\midrule
\multirow{6}{*}{attack\_tree\_} & overall\_quality & Holistic assessment of the attack tree \\
& structural\_quality & Depth, branching, organization \\
& technical\_realism & Feasibility of attack techniques \\
& attack\_path\_logic & Logical progression from access to impact \\
& completeness & Coverage of attack vectors and phases \\
& actionability & Usefulness for defenders \\
\midrule
ttp\_ & mapping\_quality & Quality of MITRE ATT\&CK technique mapping$^*$ \\
\midrule
\multirow{4}{*}{mitigation\_} & actionability & Concrete, implementable steps \\
& specificity & Tailored to identified threat and tech stack \\
& coverage & Addresses all identified attack paths \\
& technical\_accuracy & Correctness of recommended controls \\
\bottomrule
\multicolumn{3}{l}{\footnotesize $^*$Uses specialized scale: excellent, good, acceptable, poor, no\_mapping.}
\end{tabular}
\end{table}

\section{Cosine Similarity for TTP Retrieval}
\label{app:cosine}

Given an attack step description $d_i$ and ATT\&CK technique description $t_j$, ATTACK-BERT produces embedding vectors $\mathbf{e}(d_i), \mathbf{e}(t_j) \in \mathbb{R}^{768}$. The cosine similarity is:

\begin{equation}
\text{sim}(d_i, t_j) = \frac{\sum_{k=1}^{768} e_k(d_i) \cdot e_k(t_j)}{\sqrt{\sum_{k=1}^{768} e_k(d_i)^2} \cdot \sqrt{\sum_{k=1}^{768} e_k(t_j)^2}}
\end{equation}

The top-$K$ candidates are retrieved by sorting all $M$ techniques by descending similarity and selecting the first $K$ entries above threshold $\tau$:

\begin{equation}
\text{candidates}(d_i) = \{t_j \mid \text{sim}(d_i, t_j) \geq \tau\} \text{ sorted by } \text{sim}(d_i, t_j) \text{ desc, truncated to } K
\end{equation}

In our experiments, $K=5$ and $\tau=0.3$.
